% ============================================================================
% s2rcod_style.tex -- THE single style file for the S2R-COD figure set.
%
% \input THIS IN THE PREAMBLE, after \usepackage{tikz}.
% Every figure body (fig_method, fig_results_block, exp/fig_*) uses only the
% styles defined here and defines ZERO styles of its own. A \definecolor,
% \tikzset or \newcommand anywhere else in figures/ is a defect.
%
% Design rules enforced here:
%   * COLOUR IS NEVER THE ONLY CHANNEL. Every distinction is carried first by
%     shape or line style, and colour only reinforces it. The figures remain
%     readable in greyscale and to a colour-blind reader. This preserves the
%     discipline already declared in figures/methodology.tex:8-9.
%   * Palette is Okabe-Ito (colour-blind safe), used as dark strokes over very
%     light fills so that black \scriptsize text stays legible on every node.
%   * No gradients, no blur, no drop shadows, no bitmaps.
%
% Two independent marking axes (see figures/METHOD_SPEC.md sections 0 and 8):
%   SCOPE      -- from rebuild/TIER_SEGREGATION.md section 1. Written as a word
%                 (gen / signal+data / indep), NEVER as a bare tier digit,
%                 because T1/T2/T3 is overloaded four ways in this repository
%                 (rebuild/FINAL_RESULTS.md:42-53).
%   CONFIDENCE -- solid fill   = has an EXP block AND a pre-registered threshold
%                 hatched fill = has an EXP block, but the reading is post-hoc
%                 dashed border= no EXP block  (barred from both main figures)
% ============================================================================
%
% NOTE: macro names are letters-only (\scod...), because a TeX control sequence
% cannot contain a digit -- \s2foo tokenizes as \s followed by the text '2foo'.
% The tikzset KEY names keep the s2 prefix; keys are not control sequences.
% ============================================================================

\usetikzlibrary{arrows.meta,positioning,calc,fit,backgrounds,shapes.geometric,
                shapes.multipart,shapes.misc,chains,matrix,
                decorations.pathreplacing,decorations.pathmorphing,
                patterns,patterns.meta,quotes,intersections}

% ---------------------------------------------------------------- palette ---
% Okabe-Ito base hues (strokes).
\definecolor{s2blue}{HTML}{0072B2}   % data objects
\definecolor{s2green}{HTML}{009E73}  % trained modules
\definecolor{s2purple}{HTML}{CC79A7} % frozen / offline modules
\definecolor{s2orange}{HTML}{D55E00} % THE accent: feedback + selected path
\definecolor{s2amber}{HTML}{E69F00}  % decisions / gates
\definecolor{s2grey}{HTML}{595959}   % deterministic operators, rules, frames
\definecolor{s2lightgrey}{HTML}{8C8C8C}
% Very light fills (kept above 88% lightness so black text stays legible).
\definecolor{s2bluefill}{HTML}{DEEDF7}
\definecolor{s2greenfill}{HTML}{DCF0EA}
\definecolor{s2purplefill}{HTML}{F7E6F0}
\definecolor{s2orangefill}{HTML}{FBE5D8}
\definecolor{s2amberfill}{HTML}{FCF1DA}
\definecolor{s2greyfill}{HTML}{EFEFEF}
\definecolor{s2zonefill}{HTML}{F7F7F7}

% ------------------------------------------------------------ base metrics --
\tikzset{
  s2base/.style={
    line width=0.6pt, rounded corners=2pt, align=center,
    font=\scriptsize, inner sep=2.4pt, minimum height=6.6mm,
    % Hyphenation off: node text is short and ragged reads better than "ev-ery".
    % execute at begin node runs INSIDE the node box, which is where the
    % penalties must be set; /utils/exec would run during option parsing.
    execute at begin node={\hyphenpenalty=10000\exhyphenpenalty=10000},
  },
  % ---- NODE CLASSES.  Shape and border carry the meaning; colour reinforces.
  % data object: ROUNDED rectangle, thin border.
  s2data/.style   ={s2base, draw=s2blue,   fill=s2bluefill,   text width=20mm},
  % deterministic operator: SHARP rectangle.
  s2op/.style     ={s2base, draw=s2grey,   fill=s2greyfill,   text width=20mm,
                    sharp corners},
  % a deterministic operator on the SELECTED path (the arm under test).
  % Accent border only -- the shape is unchanged, so the distinction survives
  % greyscale as a line-weight difference rather than a hue difference.
  s2opsel/.style  ={s2base, draw=s2orange, fill=s2greyfill, text width=20mm,
                    sharp corners, line width=1.1pt},
  % trained module: SHARP rectangle, THICK border.
  s2trained/.style={s2base, draw=s2green,  fill=s2greenfill,  text width=22mm,
                    sharp corners, line width=1.1pt},
  % frozen module: SHARP rectangle, DOUBLE border.
  s2frozen/.style ={s2base, draw=s2purple, fill=s2purplefill, text width=22mm,
                    sharp corners, double, double distance=0.9pt},
  % gate / decision: CHAMFERED rectangle.
  s2gate/.style   ={s2base, draw=s2amber!85!black, fill=s2amberfill,
                    text width=20mm, chamfered rectangle,
                    chamfered rectangle xsep=3pt, chamfered rectangle corners={north west,south east}},
  % outcome: rounded rectangle, thick border, white ground.
  s2outcome/.style={s2base, draw=s2grey, fill=white, text width=24mm,
                    line width=1.0pt},
  % left-aligned prose box (MEASUREMENT zone, notes). Ragged right, no hyphens.
  s2textbox/.style={s2base, draw=s2grey, fill=s2greyfill, sharp corners,
                    align=left},
  % small annotation, no frame
  s2note/.style   ={font=\scriptsize, align=center, inner sep=1.2pt, text=s2grey,
                    execute at begin node={\hyphenpenalty=10000\exhyphenpenalty=10000}},
  s2tiny/.style   ={font=\tiny, align=center, inner sep=1.0pt, text=s2grey,
                    execute at begin node={\hyphenpenalty=10000\exhyphenpenalty=10000}},
  % monospace hyperparameter annotation (brief: "small monospace annotation")
  s2hyper/.style  ={font=\ttfamily\tiny, align=center, inner sep=1.2pt,
                    text=s2grey},
%
  % ---- CONFIDENCE OVERLAYS (applied on top of a node class) ----------------
  % pre-registered: the default, nothing added.
  s2prereg/.style ={},
  % post-hoc reading: hatched. Hatching is a texture, so it survives greyscale.
  s2posthoc/.style={pattern={Lines[angle=45,distance=1.5pt,line width=0.28pt]},
                    pattern color=s2lightgrey},
  % no EXP block: dashed border. Barred from fig:method and fig:results.
  s2noexp/.style  ={dashed, draw=s2lightgrey, fill=white},
  % A path or quantity that EXISTS but the pipeline cannot use at the point of
  % decision (e.g. an endpoint-measured signal at allocation time). Distinct
  % from s2unresolved, which is reserved for |Delta| <= 2*sigma_hat.
  s2unavailable/.style={draw=s2lightgrey, fill=white, dotted, text=s2grey},
  % a quantity that is NOT resolved at the stated bar -- never drawn as a win.
  s2unresolved/.style={draw=s2grey, fill=s2greyfill, dash pattern=on 2pt off 1.2pt},
%
  % ---- EDGE CLASSES --------------------------------------------------------
  s2fwd/.style      ={draw=s2grey, line width=0.6pt,
                      -{Stealth[length=1.8mm,width=1.3mm]}},
  % THE contribution edge: dashed + accent + thicker. Dash carries it alone.
  s2feedback/.style ={draw=s2orange, line width=0.9pt, dashed,
                      -{Stealth[length=1.8mm,width=1.3mm]}},
  s2selected/.style ={draw=s2orange, line width=1.0pt,
                      -{Stealth[length=1.8mm,width=1.3mm]}},
  s2rejected/.style ={draw=s2lightgrey, line width=0.5pt, dotted,
                      -{Stealth[length=1.6mm,width=1.1mm]}},
  % an offline / out-of-loop path: long dash, no accent.
  s2offline/.style  ={draw=s2purple, line width=0.6pt, dash pattern=on 3pt off 1.6pt,
                      -{Stealth[length=1.8mm,width=1.3mm]}},
  s2edgelabel/.style={font=\tiny, text=s2grey, inner sep=1.4pt,
                      fill=white, fill opacity=0.85, text opacity=1},
%
  % ---- GROUPING ------------------------------------------------------------
  s2group/.style    ={draw=s2lightgrey, line width=0.5pt, dash pattern=on 1.6pt off 1.4pt,
                      rounded corners=2.5pt, inner sep=3.2pt, fill=s2zonefill},
  s2grouplabel/.style={font=\tiny\bfseries, text=s2grey, anchor=north west,
                      inner sep=1.4pt},
%
  % ---- EXPERIMENT CARD (the four-zone per-experiment template) -------------
  % Zone 1 QUESTION, zone 2 SETUP, zone 3 MEASUREMENT, zone 4 OUTCOME.
  % fill=none is deliberate: the frame is drawn LAST so its border sits on top
  % of the content, and any fill would paint over the whole card.
  s2card/.style     ={draw=s2grey, line width=0.7pt, rounded corners=2.5pt,
                      inner sep=0pt, fill=none},
  s2question/.style ={font=\scriptsize\itshape, align=left, text=black,
                      inner sep=2.4pt, anchor=north west,
                      execute at begin node={\hyphenpenalty=10000\exhyphenpenalty=10000}},
  s2zonelabel/.style={font=\tiny\bfseries, text=s2grey, anchor=north west,
                      inner sep=0pt},
  s2rule/.style     ={draw=s2lightgrey, line width=0.4pt},
  % outcome banners. REFUTED/NULL must never look like a success.
  s2verdictnull/.style={s2base, draw=s2grey, line width=1.0pt, fill=s2greyfill,
                      font=\scriptsize\bfseries, text=black, sharp corners},
  s2verdictpos/.style ={s2base, draw=s2green, line width=1.0pt, fill=s2greenfill,
                      font=\scriptsize\bfseries, text=black, sharp corners},
  s2verdictref/.style ={s2base, draw=s2orange, line width=1.0pt, fill=s2orangefill,
                      font=\scriptsize\bfseries, text=black, sharp corners},
  % A finding that CONSTRAINS the method. Deliberately not green: a scoping
  % limit must never be drawn in the colour a success would use.
  s2verdictcon/.style ={s2base, draw=s2grey, line width=1.0pt, fill=s2greyfill,
                      font=\scriptsize\bfseries, text=black, sharp corners},
}

% ------------------------------------------------------------- text macros --
% Scope badge. Always a WORD, never a bare tier digit.
\newcommand{\scodscope}[1]{{\tiny\sffamily\bfseries\textcolor{s2grey}{[#1]}}}
\newcommand{\scodgen}{\scodscope{scope:\,generator}}
\newcommand{\scodsig}{\scodscope{scope:\,signal+data}}
\newcommand{\scodind}{\scodscope{scope:\,method-independent}}
% Seed-count marker. A single-seed number must be visibly single-seed.
\newcommand{\scodseeds}[1]{{\tiny\textcolor{s2grey}{$n{=}#1$ seeds}}}
\newcommand{\scodsingle}{{\tiny\bfseries\textcolor{s2orange}{single seed}}}
% Metric names, spelled exactly as Eval/metrics.py computes them.
\newcommand{\scodSa}{$S_\alpha$}
\newcommand{\scodMAE}{\textsc{mae}}
\newcommand{\scodsigmahat}{$\hat{\sigma}$}
\newcommand{\scodbar}{$2\hat{\sigma}$}
\newcommand{\scodrho}{$\rho$}
\newcommand{\scoddelta}[1]{$\Delta(#1)$}
% "within noise" / "not resolved" -- the phrase the repo uses, used verbatim.
\newcommand{\scodwithinnoise}{\textsc{within noise}}

% ------------------------------------------------------------------ legend --
% \scodlegend{<anchor coordinate>} -- the standard legend block, top-left
% anchored at <anchor coordinate>. ONE macro, so the legend is identical in
% content, order and geometry in every figure that carries it.
% Layout: 4 columns x 3 rows, column pitch 34mm, row pitch 5mm. The last
% label must END inside 139mm = 394.02bp, which is < the ICLR \textwidth of
% 5.5in = 396bp (iclr2027_conference.sty:49). Verified by pdfinfo in Phase 3.
% Positions are explicit offsets from the anchor, never a relative chain, so
% the block cannot drift onto the diagram.
\newcommand{\scodlegitem}[5]{% #1=xmm #2=ymm #3=style #4=anchor #5=label
  % anchor=west, not the default centre: a centred glyph would hang half its
  % width to the LEFT of the anchor and push the picture's bounding box past
  % the ICLR \textwidth. Verified with a current-bounding-box probe.
  \node[#3, text width=0pt, minimum width=5.2mm, minimum height=3.0mm,
        inner sep=0pt, rounded corners=1pt, anchor=west]
       (scodlg#1#2) at ($(#4)+(#1mm,#2mm)$) {};
  \node[s2tiny, anchor=west] at ($(scodlg#1#2.east)+(1.2mm,0)$) {#5};}
\newcommand{\scodlegedge}[5]{% #1=xmm #2=ymm #3=style #4=anchor #5=label
  \draw[#3] ($(#4)+(#1mm,#2mm)$) -- ++(5.2mm,0);
  \node[s2tiny, anchor=west] at ($(#4)+(#1mm+6.4mm,#2mm)$) {#5};}
\newcommand{\scodlegend}[1]{%
  \begin{scope}
    % row 1 -- node classes (shape and border carry the distinction)
    \scodlegitem{0}{0}{s2data}{#1}{data object}
    \scodlegitem{34}{0}{s2op}{#1}{deterministic op.}
    \scodlegitem{68}{0}{s2trained}{#1}{trained}
    \scodlegitem{102}{0}{s2frozen}{#1}{frozen / offline}
    % row 2 -- confidence axis + the uncertainty convention
    \scodlegitem{0}{-5}{s2outcome}{#1}{pre-registered}
    \scodlegitem{34}{-5}{s2outcome, s2posthoc}{#1}{post-hoc reading}
    \scodlegitem{68}{-5}{s2outcome, s2unresolved}{#1}{$|\Delta|\le2\hat{\sigma}$: unresolved}
    \scodlegitem{102}{-5}{s2gate}{#1}{gate / decision}
    % row 3 -- edge classes
    \scodlegedge{0}{-10}{s2fwd}{#1}{forward}
    \scodlegedge{34}{-10}{s2feedback}{#1}{feedback (contribution)}
    \scodlegedge{68}{-10}{s2offline}{#1}{offline, out of loop}
    \scodlegedge{102}{-10}{s2rejected}{#1}{rejected path}
  \end{scope}}

% ------------------------------------------------- per-experiment card head --
% \scodcardhead{<width mm>}{<EXP id>}{<scope macro>}{<question>}
% Drawn with the card's TOP-LEFT corner at (0,0) in a picture using x=1mm,y=1mm.
% The question box has a FIXED height, so the question line and the rule beneath
% it sit at the same y in every card -- which is what makes the library read as
% one system. Keep the question to at most two lines.
% Zone geometry, shared by every per-experiment figure:
%     y =   0      card top
%     y =  -4.6    EXP id / scope row baseline
%     y =  -6.0    question box top
%     y = -15.0    rule under the question   (\scodcardqbot)
%     y = -17.0    SETUP / MEASUREMENT zone top
%   outcome banner is placed by each figure and closed with \scodcardframe
\newcommand{\scodcardqbot}{-15}
\newcommand{\scodcardhead}[4]{%
  \node[font=\scriptsize\bfseries, anchor=north west, inner sep=0pt]
       at (1.8,-1.2) {#2};
  \node[anchor=north east, inner sep=0pt] at (#1-1.8,-1.2) {#3};
  \draw[s2rule] (0,-5.4) -- (#1,-5.4);
  \node[s2question, text width=#1mm-4mm, minimum height=8mm, anchor=north west]
       at (1.8,-6.0) {#4};
  \draw[s2rule] (0,\scodcardqbot) -- (#1,\scodcardqbot);
  \node[s2zonelabel, anchor=north west] at (1.8,\scodcardqbot-0.6) {SETUP};}
% \scodcardheadp{<width mm>}{<EXP id>}{<scope macro>}{<question>}
% Same head, but WITHOUT the SETUP zone label -- for a MERGED figure whose
% body is split into labelled panels (a)/(b)/(c) instead of SETUP/MEASUREMENT.
% The question line and both rules sit at identical y, so a paneled card still
% lines up with a four-zone card in a contact sheet.
\newcommand{\scodcardheadp}[4]{%
  \node[font=\scriptsize\bfseries, anchor=north west, inner sep=0pt]
       at (1.8,-1.2) {#2};
  \node[anchor=north east, inner sep=0pt] at (#1-1.8,-1.2) {#3};
  \draw[s2rule] (0,-5.4) -- (#1,-5.4);
  \node[s2question, text width=#1mm-4mm, minimum height=8mm, anchor=north west]
       at (1.8,-6.0) {#4};
  \draw[s2rule] (0,\scodcardqbot) -- (#1,\scodcardqbot);}
% \scodcardframe{<width mm>}{<bottom y mm>} -- the outer card border, drawn last.
\newcommand{\scodcardframe}[2]{%
  \draw[s2card] (0,0) rectangle (#1,#2);}
% \scodmeaslabel{<x mm>} -- the MEASUREMENT zone label, aligned with SETUP.
\newcommand{\scodmeaslabel}[1]{%
  \node[s2zonelabel, anchor=north west] at (#1,\scodcardqbot-0.6) {MEASUREMENT};}
